% Setting up the document class and essential packages
\documentclass[a4paper,12pt]{article}
\usepackage{amsmath, amsthm, amssymb,color,braket}
\usepackage{geometry}
\geometry{margin=0.5in} % Set 1.5-inch margins on all sides

\usepackage{parskip}
\renewcommand{\familydefault}{\sfdefault}


\newtheoremstyle{mystyle}%                % Name
{}%                                     % Space above
{}%                                     % Space below
{\itshape}%                             % Body font
{}%                                     % Indent amount
{\bfseries}%                            % Theorem head font
{.}%                                    % Punctuation after theorem head
{ }%                                    % Space after theorem head, ' ', or \newline
{}%     


\theoremstyle{mystyle}
% Configuring theorem-like environments
\newtheorem{theorem}{Theorem}[section]
\newtheorem{lemma}[theorem]{Lemma}
\newtheorem{corollary}[theorem]{Corollary}

\theoremstyle{mystyle}
\newtheorem{proposition}{Proposition}[section]
\newtheorem{definition}{\color{blue} Definition}[section]
\newtheorem{remark}{\color{red} Remark}[section]
\newtheorem{example}{Example}[section]

% Setting up section numbering
\numberwithin{equation}{section}
\title{Second Quantization}
\begin{document}
	\maketitle
	We consider the second Quantization process of identical particle.Assume the single particle Hilbert space is denoted as $\mathcal{H}$,then for $n$ particles,the Hilbert space is (by basic principles of QM)
	$$\mathcal{H}_n=\mathcal{H}\otimes \mathcal{H}\otimes  \cdots \otimes \mathcal{H}$$
	We want a system where the number of particles is not fixed.So we have to collect all Hilbert spaces with different particle number.Define the many-body Hilbert space as
	$$\mathcal{H}_\infty=\mathcal{H}_0\oplus \mathcal{H}_1\oplus\cdots =\oplus_{n=0}^\infty \mathcal{H}_n$$
\end{document}